\documentclass[11pt]{article}

\usepackage[margin=1in]{geometry}
\usepackage{amsmath,amssymb}
\usepackage{algorithm}
\usepackage{algpseudocode}
\usepackage{booktabs}
\usepackage{graphicx}
\usepackage{xcolor}
\usepackage{caption}
\usepackage{float}
\usepackage{enumitem}
\usepackage[colorlinks=true,linkcolor=blue,citecolor=blue,urlcolor=blue]{hyperref}
\usepackage{listings}
\usepackage{longtable}
\usepackage{microtype}
\usepackage{makecell}
\setlength{\emergencystretch}{2em}

\definecolor{codebg}{RGB}{245,245,245}
\lstset{
  basicstyle=\ttfamily\footnotesize,
  backgroundcolor=\color{codebg},
  breaklines=true,
  frame=single,
  columns=fullflexible,
  captionpos=b
}

\newcommand{\pendhrl}{Pen-DHRL-TD}
\newcommand{\code}[1]{\texttt{#1}}

\title{\pendhrl{}: Teacher-Distilled Deep Hierarchical\\
Reinforcement Learning for Autonomous Penetration\\
Testing via Locally-Hosted LLM Subgoal Supervision}

\author{NASimEmu-enhanced Project}
\date{\today}

\begin{document}
\maketitle

\begin{abstract}
We present \pendhrl{}, a deep hierarchical reinforcement learning (DHRL) agent for
autonomous network penetration testing in the NASimEmu simulator, augmented with
offline supervision from a locally-hosted large language model (LLM) acting as a
\emph{teacher} over the agent's high-level subgoal space. The agent (\code{NASimNetDHRL})
factorizes its policy into a recurrent graph-neural-network \emph{manager} that
selects one of eight named semantic subgoals (e.g.\ \code{GAIN\_INITIAL\_ACCESS},
\code{ESCALATE\_PRIVILEGE}) and persists it for a learned horizon, and a
goal-conditioned graph-attention \emph{worker} that emits low-level scan/exploit/pivot
actions. A locally-hosted, free, open-weight model (\code{qwen2.5:3b-instruct}, served
via Ollama) is queried \emph{offline only} -- never inside the training loop's forward
pass -- against sanitized, partially-observable state summaries collected by a
triggered-sampling data-collection pass, producing a provenance-tracked dataset of
(state, goal, confidence) triples. This dataset drives a four-stage training schedule:
(1) supervised warm-start of the subgoal head, (2--3) joint PPO training with an
exponentially-decaying confidence-weighted distillation loss, and (4) a
guaranteed reinforcement-learning-only tail. We report the full architecture,
teacher pipeline, training algorithm, and a live explainability mechanism that
exposes the manager's subgoal-selection histogram and distillation diagnostics every
epoch. On a 12-subnet, 72--96-host dynamic corporate scenario family, we report
end-to-end training dynamics, a systematic class-imbalance issue discovered via
dataset audit, a checkpoint-selection pitfall uncovered during held-out evaluation,
and cross-scenario generalization results against both a metric-frozen checkpoint of
the same run and an independently trained, non-distilled Pen-DHRL baseline of
identical architecture and parameter count, on which the distilled agent comes out
ahead on every tested scenario and every reported metric.
\end{abstract}

\tableofcontents
\newpage

% =====================================================================
\section{Introduction}
% =====================================================================

Autonomous penetration testing is naturally hierarchical: a human operator does not
reason step-by-step over individual port scans and exploit attempts, but over
higher-level intentions -- ``enumerate this host'', ``escalate on that foothold'',
``reduce detection risk before proceeding''. Flat reinforcement-learning policies
trained directly on primitive actions must rediscover this structure from reward
alone, which is slow and sample-inefficient on large, sparse-reward network
topologies. \pendhrl{} addresses this with an explicit two-level hierarchy: a
\emph{manager} that commits to a persistent semantic subgoal, and a \emph{worker}
that acts under that subgoal's guidance.

The manager's subgoal space is, by construction, an arbitrary learned embedding bank
with no inherent meaning -- goal index 3 is only ``escalate privilege'' because the
worker's policy gradient happens to shape it that way. This project's central idea is
to bind those otherwise-anonymous subgoal slots to a fixed, human-readable ontology
(Section~\ref{sec:ontology}) and use an LLM, queried entirely offline against
sanitized state summaries, as a cheap source of weak supervision over \emph{which}
named goal a reasonable planner would pick at a given state. This is emphatically
\textbf{not} an LLM-in-the-loop agent: the teacher is never queried during environment
interaction or during the PPO forward/backward pass. It is queried once, offline, to
build a static dataset, exactly as a human-labeled dataset would be built, and the
resulting labels are distilled into the manager's subgoal head via an auxiliary
supervised loss that runs alongside -- never inside -- the reinforcement-learning
update.

\subsection{Contributions}
\begin{itemize}[leftmargin=1.4em]
  \item A hierarchical graph-policy architecture (\code{NASimNetDHRL}) with a
        recurrent GNN manager, a learned subgoal bank, a persistence (``switch'')
        gate, and a goal-conditioned graph-attention worker (Section~\ref{sec:arch}).
  \item A fixed 8-way semantic goal ontology and an offline LLM-teacher pipeline
        (state sanitization $\to$ JSON-schema-constrained local LLM query $\to$
        anti-hallucination validation $\to$ provenance-tracked dataset) that produces
        confidence-weighted labels without ever exposing hidden simulator state to
        the teacher (Sections~\ref{sec:ontology}--\ref{sec:teacher}).
  \item A triggered-sampling data-collection procedure that queries the teacher only
        at strategically informative moments (episode start, new subnet, first
        compromise, privilege change, IDS alert, repeated failure, periodic fallback,
        high manager entropy) under a diversified mixture of rollout policies
        (Section~\ref{sec:collection}).
  \item A four-stage training schedule combining supervised warm-start, decaying
        joint distillation, and a guaranteed RL-only tail, plus an optional
        potential-based reward-shaping channel over the same semantic goal space
        (Sections~\ref{sec:distill}--\ref{sec:training}).
  \item A live, per-epoch \emph{explainability} mechanism that prints the manager's
        realized subgoal-selection distribution and distillation diagnostics in
        human-readable form during training -- the mechanism that, in this project's
        own training run, first exposed a subgoal-selection collapse later confirmed
        through held-out evaluation (Section~\ref{sec:explainability}).
  \item An end-to-end empirical study on a 100-host dynamic corporate scenario family:
        dataset audit statistics, training dynamics, a checkpoint-selection pitfall,
        and cross-scenario generalization results against both a same-run checkpoint
        and an independently trained, non-distilled baseline (Section~\ref{sec:results}).
\end{itemize}

% =====================================================================
\section{Environment and Task}
\label{sec:env}
% =====================================================================

The agent operates in NASimEmu, a partially-observable network-attack-simulation
gym environment. An episode consists of a randomly-instantiated corporate network;
the agent observes only what it has actively discovered (host existence, OS,
services, processes, access level, IDS detection level) and must chain
scan/exploit/privilege-escalation actions to compromise hosts, subject to a fixed
per-episode step budget.

\subsection{Scenario family}
All experiments use the \code{corp\_100hosts\_dynamic.v2} scenario family: 12
labeled subnets (DMZ, user, service, db, dev, app $\times$3, web, office, research,
backups) connected in a fixed topology matrix, each instantiated at reset with a
random host count in $[6,8]$ (72--96 hosts total per episode). Per-subnet
sensitive-host probability ranges from $0.0$ (DMZ/user) to $1.0$ (db), giving a
non-uniform, subnet-dependent value landscape. A 32-entry extended service catalog
and a process catalog (each tagged with an OS) are drawn per host; three services
(\code{mysql}, \code{postgres}, \code{mssql}) are up-weighted on sensitive hosts.
Five concrete scenario instances are used across this study -- \code{base},
\code{varA}, \code{varB}, \code{test}, \code{bridge} -- which share topology and
service/process catalogs but differ in per-exploit/privesc \emph{success
probability} and \emph{cost}, giving a controllable difficulty gradient
(Table~\ref{tab:scenario-difficulty}).

\begin{table}[H]
\centering
\caption{Mean exploit success probability across the 8-entry exploit catalog, per
scenario instance (lower = harder). Costs increase in the same order.}
\label{tab:scenario-difficulty}
\begin{tabular}{lcc}
\toprule
Scenario & Role & Mean exploit prob.\\
\midrule
\code{corp\_100hosts\_dynamic} (base) & training & $\approx 0.79$ \\
\code{corp\_100hosts\_dynamic\_varA} & training & $\approx 0.63$ \\
\code{corp\_100hosts\_dynamic\_varB} & training & $\approx 0.58$ \\
\code{corp\_100hosts\_dynamic\_bridge} & held-out eval & $\approx 0.66$ \\
\code{corp\_100hosts\_dynamic\_test} & held-out eval & $\approx 0.53$ (hardest) \\
\bottomrule
\end{tabular}
\end{table}

\subsection{Observation and action space}
The agent consumes a \code{graph\_v2} observation: nodes are either subnets or hosts
(a leading type bit distinguishes them), host nodes carry an 80-dimensional feature
vector (discovery/reachability/compromise/access flags, per-service and per-process
booleans, IDS detection level and threshold, sensitivity-irrelevant fields only --
see Section~\ref{sec:teacher} for the same partial-observability constraint applied
to the LLM teacher), and an 8-dimensional sinusoidal positional encoding is appended
per node. Edges follow host reachability. The action space is a (target node,
discrete action id) pair over 18 action types (4 fixed scan actions plus the
scenario's exploit/privilege-escalation catalog), masked so that subnet nodes and
graphs with no valid host targets cannot be selected. An explicit terminate action
(\code{-use\_a\_t}) lets the agent end an episode early rather than exhausting the
400-step budget.

% =====================================================================
\section{Hierarchical Policy Architecture}
\label{sec:arch}
% =====================================================================

\code{NASimNetDHRL} (212{,}640 parameters at the hyperparameters used in this study)
is a manager--worker hierarchical graph policy, structurally inspired by FeudalNet-
and GNN-LSTM-style architectures adapted to a variable-size graph observation.

\subsection{Manager}
Node features (with positional encoding) are linearly embedded
($\mathbb{R}^{88}\!\to\!\mathbb{R}^{64}$) and passed through a 3-iteration
multi-message-passing GNN with a recurrent global ``last-GRU'' readout
(\code{Multi\-Message\-Passing\-With\-Global\-Node\-And\-Last\-GRU\-Global}), producing per-node
embeddings $x_{\text{mgr}}$ and a single recurrent global summary $x_{\text{global}}
\in \mathbb{R}^{64}$ per graph (episode). $x_{\text{global}}$ is layer-normalized and
additively biased by an IDS-aware projection of the graph's mean detection
level/threshold/multiplier (\code{ids\_projection}), mirroring a FeudalGTM-style
stealth bias: the manager's context is nudged by how ``hot'' the network currently
looks, encouraging (without hard-constraining) detection-averse planning.

\subsection{Subgoal selection and persistence}
A subgoal head $\pi_{\text{goal}}(\cdot\mid x_{\text{global}})$
(2-layer MLP $\to$ 8 logits) defines a categorical distribution over the
$K{=}8$ subgoal slots. A separate \emph{switch} gate
$\sigma_{\text{switch}}(x_{\text{global}}) \in (0,1)$ (sigmoid MLP) is sampled as a
Bernoulli \emph{persistence-override} signal each step: the manager only re-samples a
new subgoal when either (a) the previous subgoal's horizon counter has reached zero,
or (b) the switch gate fires. This is what makes subgoals \emph{temporally extended}
rather than re-selected every step -- purely reactive per-step goal selection would
give the worker no stable target to condition a multi-step plan on. Once selected,
subgoal index $k$ indexes a learned embedding \code{goal\_bank}$[k] \in
\mathbb{R}^{128}$, refined by concatenation with $x_{\text{global}}$ through a small
MLP (\code{goal\_refine}) to produce the state-conditioned goal vector actually
broadcast to the worker. The joint action probability (used for PPO's importance
ratio) correctly folds in the probability of \emph{whichever} of these three
branches produced the current subgoal (forced switch, policy-initiated switch, or
persistence) -- Section~\ref{sec:arch}'s Algorithm~\ref{alg:forward}, lines
constructing \code{subgoal\_factor}, give the exact three-way case split.

\subsection{Worker}
The refined goal vector is broadcast to every node in its graph and concatenated
with the manager's per-node embedding; a 2-head graph-attention layer
(\code{GATConv}) followed by layer norm produces goal-conditioned node features,
linearly projected to per-node action logits. Illegal targets (subnet nodes, or
every node in a graph with no valid host target) are masked to $-\infty$ before a
per-graph softmax, and the realized (target, action) pair is sampled by a segmented
categorical sample over the flattened, graph-grouped logits.

\subsection{Value, termination, and time pressure}
A shared context vector (global embedding $\oplus$ goal vector $\oplus$ normalized
elapsed-time fraction) feeds both a value head and a termination head (sigmoid). An
explicit linear-plus-threshold time penalty is subtracted from the value estimate as
the episode approaches its step budget, and termination is additionally gated to
require both a minimum elapsed fraction (25\% of the step limit) and a non-positive
value estimate -- this two-sided gate (Algorithm~\ref{alg:forward}, termination
block) is what prevents the policy from learning to terminate near-instantly for a
cheap zero-cost episode, while still allowing genuinely early, successful
termination once real progress has been made.

\begin{algorithm}[H]
\caption{\code{NASimNetDHRL.forward} -- Manager--Worker Hierarchical Policy Step}
\label{alg:forward}
\begin{algorithmic}[1]
\Require batch of graph observations $\{s_i\}$, persistent state
  $(g^{\text{idx}}_{t-1}, \tau_{t-1}, h_{t-1})$ (subgoal index, steps-remaining, GRU
  hidden state) carried from the previous call
\Ensure actions $\{(v_i, a_i)\}$, value $V$, joint action probability $\pi$, raw
  actions $(a^{\text{sel}}, \text{term}, g^{\text{idx}}_t)$
\State $x \gets \text{Embed}(\text{concat}(\text{node\_feats}, \text{PosEnc}))$
\State $x_{\text{mgr}}, x_{\text{global}} \gets \text{ManagerGNN}(x, \text{edges}, h_{t-1})$
  \Comment{recurrent global readout}
\State $x_{\text{global}} \gets \text{LayerNorm}(x_{\text{global}}) + \text{IDSProj}(\bar\phi_{\text{ids}})$
\State $\ell_{\text{goal}} \gets \text{SubgoalHead}(x_{\text{global}})$;\quad
       $p_{\text{goal}} \gets \text{softmax}(\ell_{\text{goal}})$
\State $\sigma \gets \text{clamp}(\text{SubgoalSwitch}(x_{\text{global}}), \epsilon, 1-\epsilon)$
\If{$\tau_{t-1} \le 0$} \Comment{horizon exhausted: forced re-sample}
  \State $\text{need\_new} \gets \text{true}$
\EndIf
\State $\text{switch} \gets \text{need\_new} \lor \text{Bernoulli}(\sigma)$ \Comment{sample-mode; argmax test-mode}
\If{$\text{switch}$}
  \State $g^{\text{idx}}_t \sim \text{Categorical}(p_{\text{goal}})$ \Comment{argmax at eval time}
  \State $\tau_t \gets \text{goal\_horizon}$ \Comment{reset persistence counter (=4)}
\Else
  \State $g^{\text{idx}}_t \gets g^{\text{idx}}_{t-1}$;\quad $\tau_t \gets \tau_{t-1} - 1$
\EndIf
\State $\pi_{\text{subgoal}} \gets$ probability of \emph{whichever} branch above fired
  \Comment{3-way case: forced / policy-switch / persistence}
\State $\vec g \gets \text{GoalRefine}(\text{concat}(\text{goal\_bank}[g^{\text{idx}}_t], x_{\text{global}}))$
\State $z \gets \text{GAT}(\text{concat}(x_{\text{mgr}}, \vec g[\text{batch\_ind}]),\ \text{edges})$
\State $\ell_{\text{action}} \gets \text{ActionHead}(\text{LayerNorm}(z))$; mask illegal targets to $-\infty$
\State $a^{\text{sel}} \sim \text{Categorical}(\text{softmax}(\ell_{\text{action}}))$ per graph \Comment{segmented sample}
\State $\text{ctx} \gets \text{concat}(x_{\text{global}}, \vec g, t/T)$
\State $V \gets \text{ValueHead}(\text{ctx}) - \text{TimePenalty}(t/T)$
\State $p_{\text{term}} \gets \text{TerminationHead}(\text{ctx})$
\State $\text{term} \sim \text{Bernoulli}(p_{\text{term}}) \ \land\ (t \ge 0.25\,T) \ \land\ (V \le 0)$
  \Comment{gated early stop}
\State $\pi \gets \begin{cases} p_{\text{term}} & \text{term} \\
  \pi_{\text{action}} \cdot (1-p_{\text{term}}) \cdot \pi_{\text{subgoal}} & \text{otherwise}\end{cases}$
\State \Return $\{(v_i, a_i)\}$, $V$, $\pi$, $(a^{\text{sel}}, \text{term}, g^{\text{idx}}_t)$
\end{algorithmic}
\end{algorithm}

% =====================================================================
\section{Semantic Goal Ontology}
\label{sec:ontology}
% =====================================================================

The manager's $K{=}8$ subgoal-bank slots have no inherent meaning to the
architecture -- they are exactly as trainable as an unlabeled latent-goal DHRL
baseline. \pendhrl{} fixes a versioned index$\to$name binding
(Table~\ref{tab:ontology}) so that (a) an LLM teacher can be prompted against stable
names instead of opaque integers, (b) potential-based reward shaping functions can be
hand- or LLM-authored per named goal, and (c) every trained checkpoint stores its
ontology version and refuses to load under a mismatched version (\code{net.py}'s
\code{Net.save}/\code{load}), preventing a silent semantic mismatch between a
checkpoint's learned goal-bank and a later ontology revision.

\begin{table}[H]
\centering
\caption{Fixed semantic goal ontology (version 1).}
\label{tab:ontology}
\begin{tabular}{clp{7.5cm}}
\toprule
Idx & Name & Description (as given to the LLM teacher) \\
\midrule
0 & \code{DISCOVER\_SUBNET} & Reveal new reachable subnet structure \\
1 & \code{ENUMERATE\_HOST} & Learn OS/services/processes of a discovered host \\
2 & \code{GAIN\_INITIAL\_ACCESS} & Obtain user or root access on a discovered, reachable host remotely \\
3 & \code{ESCALATE\_PRIVILEGE} & Upgrade an already-held user access to root on the same host \\
4 & \code{PIVOT} & Move toward/through a newly reachable, deeper subnet \\
5 & \code{CAPTURE\_SENSITIVE\_HOST} & Prioritize progress on a host that looks high-value \\
6 & \code{REDUCE\_DETECTION} & Prefer a lower-detection-risk action or target over a riskier one \\
7 & \code{RECOVER\_OR\_REPLAN} & Abandon a goal that has repeatedly failed and pick a feasible alternative \\
\bottomrule
\end{tabular}
\end{table}

% =====================================================================
\section{Offline LLM Teacher Pipeline}
\label{sec:teacher}
% =====================================================================

\subsection{Design constraints}
The teacher is queried \textbf{offline only}: never inside \code{forward()}, never
during environment interaction, never during the PPO backward pass. It is a locally
hosted, free, open-weight model (\code{qwen2.5:3b-instruct}, served via Ollama's
\code{/api/chat} at zero per-call monetary cost, JSON-schema-constrained decoding,
temperature 0) -- not a paid hosted API, since nothing in this design needs low
latency or a live network call during training. The teacher sees only a
\emph{sanitized state summary}, never the raw graph tensor or any hidden simulator
field: only hosts with \code{discovered==1} are listed, and the true
sensitive-host label (\code{HostVector.value}/\code{discovery\_value}) is never
exposed -- exactly the same restraint independently enforced in the potential-shaping
functions of Section~\ref{sec:shaping}. The teacher must choose exactly one goal from
the closed ontology of Table~\ref{tab:ontology}, optionally an already-discovered
target, and is explicitly instructed never to propose executable text (shell
commands, exploit code); a denylist-based validator enforces this on the free-text
fields regardless.

\subsection{Prompt and output contract}
The system prompt fixes the assistant's role and hard constraints; the user prompt
embeds the closed goal list with one-line descriptions and the JSON-serialized state
summary, and requests a strict-JSON response validated against a JSON schema
(\code{format} passed to Ollama) with required fields \code{goal}, \code{horizon}
($1$--$8$), \code{confidence} ($[0,1]$), \code{reason\_code}, \code{rationale}, and
optional \code{target\_subnet}/\code{target\_host}. Schema-constrained decoding
guarantees \emph{shape} only; a separate semantic validator
(\code{llm\_teacher/validator.py}) additionally rejects: any target host/subnet id
not literally present in the summary (\emph{anti-hallucination}), a target
structurally incompatible with the chosen goal (e.g.\ \code{ESCALATE\_PRIVILEGE}
against a host the agent does not already hold user access on), and any of a small
denylist of shell/code tokens appearing in the free-text fields. Up to 3 retries are
attempted on invalid/unparsable output; the final attempt is recorded with
\code{valid=False} and a \code{reject\_reason} rather than silently discarded, so
rejected records remain available for later audit.

\subsection{Optional escalation cascade}
A primary/escalation \emph{cascade} variant (\code{llm\_teacher/cascade.py}) queries
a primary model on every state and escalates to a second model only on concrete,
checkable triggers: an invalid/unparsable primary answer, a
\code{GAIN\_INITIAL\_ACCESS} goal with no reachable unclaimed host left (an
unsatisfiable goal the primary failed to notice), three consecutive failures against
the same target that the primary did not route to \code{RECOVER\_OR\_REPLAN}, or
low self-reported confidence. The escalation call uses a \emph{veto/critique} design:
the second model sees the primary's answer and either endorses it
(\code{defensible=true}) or supplies a full, independently-validated replacement --
deliberately narrower and cheaper than an independent re-decision, and not trusted
automatically just because it was consulted second. This variant is documented here
for completeness; the dataset used in Section~\ref{sec:results} was collected with
the single-model backend only (all 2{,}235 records attributed to
\code{qwen2.5:3b-instruct}).

\subsection{Heuristic teacher (ablation baseline)}
A deterministic, network-free rule-based teacher
(\code{llm\_teacher/heuristic\_teacher.py}) implements a fixed-priority decision tree
over the identical sanitized summary (escalate on any held user access
$\to$ initial-access on any enumerated-but-unclaimed host $\to$ enumerate any
reachable-but-unscanned host $\to$ capture on any root-access host $\to$ reduce
detection near an IDS threshold $\to$ replan on repeated same-target failure
$\to$ default explore) and is routed through the same validator as the LLM. Its
purpose is a controlled ablation (``if simple hand-written heuristics match the LLM's
effect, the LLM is not doing anything a much cheaper mechanism could not'') and is
not the backend used for the results in this report.

% =====================================================================
\section{Triggered Dataset Collection}
\label{sec:collection}
% =====================================================================

Labeling every visited state would be both expensive and redundant. Instead, the
teacher is queried only at \emph{strategically triggered} points, under a
deliberately diversified mixture of rollout (behavior) policies -- pure random
action selection, an untrained (randomly initialized) \code{NASimNetDHRL} at zero
training cost, or a partially/fully trained checkpoint -- so the dataset is not
narrowed to only the states a single fixed policy happens to visit (labeling only
successful expert trajectories would starve rare-but-important classes such as
\code{ESCALATE\_PRIVILEGE}).

\begin{algorithm}[H]
\caption{Triggered State--Action Labeling (Dataset Collection)}
\label{alg:collect}
\begin{algorithmic}[1]
\Require scenario, target valid-record count $N$, rollout policy $\pi_{\text{roll}}$,
  teacher backend $\mathcal{T}$, periodic interval $P$
\Ensure dataset $\mathcal{D}$ of provenance-tracked (state, goal, confidence) records
\State $n_{\text{valid}} \gets |\{r \in \mathcal{D} : r.\text{valid}\}|$
\While{$n_{\text{valid}} < N$ and total steps $<$ max\_steps}
  \State reset environment; new random $\text{episode\_id}$
  \State label state at \textbf{episode start} (always triggered); update $n_{\text{valid}}$
  \For{$t = 1, \dots, \text{step\_limit}$}
    \State $a \gets \pi_{\text{roll}}(\text{state})$; step environment; observe $(s', r, d, \text{info})$
    \State $\text{triggered} \gets \textbf{false}$
    \If{new subnet discovered} $\text{triggered} \gets \textbf{true}$ \Comment{trigger: \emph{new\_subnet\_discovered}}
    \ElsIf{new host compromised} $\text{triggered} \gets \textbf{true}$ \Comment{\emph{first\_compromise}}
    \ElsIf{access level increased on any known host} $\text{triggered} \gets \textbf{true}$ \Comment{\emph{privilege\_change}}
    \ElsIf{new IDS threshold crossing} $\text{triggered} \gets \textbf{true}$ \Comment{\emph{ids\_alert}}
    \ElsIf{last 3 actions: same target, all failed} $\text{triggered} \gets \textbf{true}$ \Comment{\emph{three\_consecutive\_failures}}
    \ElsIf{$\pi_{\text{roll}}$ is net-based \textbf{and} $H(\text{subgoal dist.}) > 0.8\,H_{\max}$} $\text{triggered} \gets \textbf{true}$
      \Comment{\emph{high\_student\_entropy}, net-based policies only}
    \ElsIf{$t \bmod P = 0$} $\text{triggered} \gets \textbf{true}$ \Comment{\emph{periodic} fallback}
    \EndIf
    \If{$d$} \textbf{break} \Comment{episode ended}
    \EndIf
    \If{$\text{triggered}$}
      \State $\text{summary} \gets \text{SanitizedSummary}(s')$ \Comment{discovered-only, no hidden fields}
      \State $\text{record} \gets \mathcal{T}(\text{summary})$ \Comment{query teacher; validate; retry $\le 3\times$}
      \State $\mathcal{D} \gets \mathcal{D} \cup \{\text{record}\}$; update $n_{\text{valid}}$ if valid
    \EndIf
  \EndFor
\EndWhile
\State \Return $\mathcal{D}$
\end{algorithmic}
\end{algorithm}

Each episode is stamped with a random \code{episode\_id}, later used by
\code{split\_dataset.py} to split train/test at the \emph{scenario-instance} level
(never placing two steps of the same episode across the boundary) via a
deterministic, seed-salted hash of \code{episode\_id} -- stable under dataset growth,
since it does not depend on shuffling the full episode-id list positionally.

% =====================================================================
\section{Confidence-Weighted Goal Distillation}
\label{sec:distill}
% =====================================================================

Only a \emph{hard-label} distillation term $\mathcal{L}_{\text{goal}}$ is implemented,
not a soft-label KL term: the teacher returns exactly one discrete goal plus a scalar
confidence, never a full probability distribution over the ontology, and fabricating
a soft target from a single discrete choice was explicitly rejected by design.

\begin{equation}
\mathcal{L}_{\text{goal}} \;=\; \frac{1}{N}\sum_{i=1}^{N} w_i \cdot \big[-\log
  \hat p_{\theta}(g_i \mid s_i)\big], \qquad
  w_i = \operatorname{clip}(c_i,\, w_{\min},\, 1), \quad w_{\min}=0.05
\label{eq:distill-loss}
\end{equation}

where $\hat p_\theta(\cdot \mid s_i) = \operatorname{softmax}(\ell_{\text{goal}}(s_i))$
is the manager's subgoal head evaluated via a side-channel forward pass
(\code{only\_subgoal\_logits=True}) that runs the manager GNN only -- skipping the
worker, value, and termination machinery entirely -- and $c_i \in [0,1]$ is the
teacher's self-reported confidence for record $i$, clipped away from zero so a
low-confidence label still contributes rather than being silently dropped.

\begin{algorithm}[H]
\caption{Confidence-Weighted Goal Distillation Loss}
\label{alg:distill}
\begin{algorithmic}[1]
\Function{GoalDistillationLoss}{$\ell_{\text{goal}} \in \mathbb{R}^{N\times 8}$, $g \in \{0,\dots,7\}^N$, $c \in [0,1]^N$}
  \State $w \gets \operatorname{clip}(c,\, 0.05,\, 1.0)$
  \State $\log \hat p \gets \log\text{-softmax}(\ell_{\text{goal}}, \text{dim}=-1)$
  \State $\text{nll}_i \gets -\log \hat p_i[g_i]$ for each $i$
  \State \Return $\text{mean}_i ( w_i \cdot \text{nll}_i )$
\EndFunction
\end{algorithmic}
\end{algorithm}

An optional random-label control (\code{--llm\_distill\_shuffle\_labels}, ablation
condition C5) permutes \code{goal\_idx} across records \emph{in place} -- preserving
the real class-frequency distribution -- so the same states and confidences train
against structurally-scrambled labels, isolating whether an observed gain comes from
real teacher knowledge or merely the regularizing effect of an extra auxiliary loss.

% =====================================================================
\section{Optional Potential-Based Reward Shaping}
\label{sec:shaping}
% =====================================================================

A separate, independently togglable channel (\code{--llm\_shaping}) applies
potential-based reward shaping directly at the reward level, over the augmented
(state, active-subgoal) space $x_t = (s_t, g_t)$:
\begin{equation}
F_t = \gamma\, \lambda_{t+1}\, \Phi(x_{t+1}) \;-\; \lambda_t\, \Phi(x_t)
\label{eq:pbrs}
\end{equation}
using the \emph{realized} next subgoal $g_{t+1}$ (never a frozen copy of $g_t$), with
terminal potentials fixed to 0 and shaping weight $\lambda$ annealed across global
training steps but held constant within a single PPO rollout window -- so
Eq.~\ref{eq:pbrs} reduces exactly to $\lambda \cdot (\gamma \Phi(x_{t+1}) -
\Phi(x_t))$ within a window, an exact instance of the corrected potential-based
shaping form (not an approximation of it), since $\lambda$'s own update granularity
is defined at the window level. Seven of the eight goals have a hand-authored,
non-constant potential $\Phi_g(\text{state}) \in [0,1]$ built only from the same
partially-observable fields the baseline network already reads (fraction of hosts
discovered, fraction scanned, fraction with user/root access, mean detection level,
etc.); \code{RECOVER\_OR\_REPLAN} is left an explicit neutral placeholder
($\Phi \equiv 0$), documented rather than silently faked. All potentials are computed
once at rollout-collection time under \code{torch.no\_grad()} and cached into the
reward \emph{before} the PPO update, never recomputed during PPO's multi-epoch
replay. \textbf{This channel was not enabled in the training run reported in
Section~\ref{sec:results}} (\code{--llm\_shaping} was not passed); it is documented
here as part of the full system design and for architectural completeness.

\begin{algorithm}[H]
\caption{Potential-Based Reward Shaping Term (per rollout window)}
\label{alg:shaping}
\begin{algorithmic}[1]
\Require rollout trace of length $T$ (=\code{ppo\_t}), each step's realized subgoal $g_\tau$, done flag $d_\tau$
\Require shaping weight $\lambda_t$ (constant within this window)
\For{$\tau = 1, \dots, T$}
  \State $\Phi_\tau \gets \Phi_{g_\tau}(s_\tau)$ \Comment{gather the active goal's potential (Sec.~\ref{sec:shaping})}
\EndFor
\For{$\tau = 1, \dots, T$}
  \If{$\tau < T$}
    \State $\Phi_{\tau+1}' \gets 0 \text{ if } d_\tau \text{ else } \Phi_{\tau+1}$
    \State $F_\tau \gets \lambda_t (\gamma \Phi_{\tau+1}' - \Phi_\tau)$
  \Else \Comment{window-boundary step: $g_{\tau+1}$ not yet known}
    \State $F_\tau \gets \lambda_t(0 - \Phi_\tau)$ if $d_\tau$ else $0$ \Comment{exact if terminal, neutral otherwise}
  \EndIf
\EndFor
\State \Return $\{F_\tau\}_{\tau=1}^T$ \Comment{added to that step's env reward before \code{net.update}}
\end{algorithmic}
\end{algorithm}

% =====================================================================
\section{Four-Stage Training Procedure}
\label{sec:training}
% =====================================================================

Training combines on-policy PPO over the environment reward with the distillation
loss of Section~\ref{sec:distill} in four stages, controlled by a single schedule
over global training step $t$ (Eq.~\ref{eq:lambda-goal-sched}):

\begin{equation}
\lambda_{\text{goal}}(t) =
\begin{cases}
  0 & \text{if } t \ge (1-f)\,T_{\text{tot}} \quad \text{(Stage 4: RL-only, forced exactly 0)} \\[2pt]
  \lambda_0 \exp(-t/\rho) & \text{otherwise \quad (Stages 2--3: fixed-then-decaying joint weight)}
\end{cases}
\label{eq:lambda-goal-sched}
\end{equation}

where $T_{\text{tot}} = {\tt max\_epochs} \times {\tt log\_rate}$ is the total
planned step budget, $f$ (flag \code{-{}-llm\-\_distill\-\_rl\-\_only\-\_frac}) is
the final fraction of training reserved as guaranteed RL-only, $\lambda_0$ (flag
\code{-{}-llm\-\_distill\-\_lambda\-\_start}, default 1.0), and $\rho$ (flag
\code{-{}-llm\-\_distill\-\_lambda\-\_rate}, default 8000 steps). The
$t \ge (1-f)T_{\text{tot}}$ branch exists because the plain exponential decay alone
only \emph{asymptotically} approaches zero and never guarantees a true RL-only phase
within a bounded run -- Stage 4 forces it exactly, once.

\begin{algorithm}[H]
\caption{Full LLM-Distilled Training Procedure}
\label{alg:train}
\begin{algorithmic}[1]
\Require dataset $\mathcal{D}$ (Algorithm~\ref{alg:collect}), warm-start epochs $E_w$,
  max epochs $E_{\max}$, schedule params $(\lambda_0, \rho, f)$
\State \textbf{Stage 1 (warm-start):} \Comment{pure supervised, $\mathcal{L}_{\text{RL}}=0$}
\For{$e = 1,\dots,E_w$}
  \For{minibatch $B \subset \mathcal{D}$ (shuffled)}
    \State $\ell_{\text{goal}} \gets \text{net}(B.\text{states}, \text{only\_subgoal\_logits}{=}\text{True})$
    \State $\mathcal{L} \gets \textsc{GoalDistillationLoss}(\ell_{\text{goal}}, B.\text{goal\_idx}, B.\text{confidence})$
    \State backprop $\mathcal{L}$; clip grad norm; optimizer step
  \EndFor
\EndFor
\State hard-sync target network to warm-started weights ($\rho_{\text{target}}{=}1.0$)
\State \textbf{Stages 2--4 (joint PPO $+$ decaying/zeroed distillation):}
\For{global step $t = 1, 2, \dots$}
  \State collect an on-policy rollout of length \code{ppo\_t} using the current policy
    (Algorithm~\ref{alg:forward} each step)
  \If{\code{--llm\_shaping} enabled}
    \State shape rollout rewards via Algorithm~\ref{alg:shaping}
  \EndIf
  \State $\mathcal{L}_{\text{RL}} \gets$ clipped-surrogate PPO loss over the rollout (Section~\ref{sec:ppo})
  \State update net via $\mathcal{L}_{\text{RL}}$ \Comment{primary gradient step}
  \State $\lambda_{\text{goal}} \gets$ Eq.~\ref{eq:lambda-goal-sched} evaluated at $t$
  \If{$\lambda_{\text{goal}} > 10^{-6}$}
    \State sample minibatch $B \subset \mathcal{D}$
    \State $\ell_{\text{goal}} \gets \text{net}(B.\text{states}, \text{only\_subgoal\_logits}{=}\text{True}, \text{reset\_hidden}{=}\text{True})$
    \State $\mathcal{L}_{\text{goal}} \gets \lambda_{\text{goal}} \cdot \textsc{GoalDistillationLoss}(\ell_{\text{goal}}, B.\text{goal\_idx}, B.\text{confidence})$
    \State backprop $\mathcal{L}_{\text{goal}}$ \textbf{as a separate gradient step} \Comment{not fused into $\mathcal{L}_{\text{RL}}$'s backward}
  \EndIf
  \State every \code{log\_rate} steps: evaluate, log explainability brief (Sec.~\ref{sec:explainability}), checkpoint
  \If{$t / \code{log\_rate} \ge E_{\max}$} \textbf{break}
  \EndIf
\EndFor
\end{algorithmic}
\end{algorithm}

The distillation gradient step is deliberately kept \emph{separate} from the PPO
update rather than fused into one combined backward pass: fusing would require
threading an auxiliary loss through the shared \code{rl.py:ppo()} update function
used by every network architecture in this codebase, which this project's scope does
not extend to. This alternating scheme optimizes the same two objectives, just not in
a single backward pass.

\subsection{PPO update}
\label{sec:ppo}
The manager's recurrence means PPO's usual multi-epoch importance-sampled replay
(\code{ppo\_k}$=3$ epochs per rollout) must replay the \emph{entire} \code{ppo\_t}-step
window in sequence from a saved hidden state $h_{s_0}$, not as an unordered
minibatch, so the manager GNN's recurrent state stays consistent with the actions
being replayed. Advantage is computed against a soft target network updated by Polyak
averaging ($\rho_{\text{target}}=0.1$), and the clipped surrogate objective
$\mathcal{L}_\pi = -\mathbb{E}\big[\min(\rho_t A_t,\, \text{clip}(\rho_t, 1{-}\epsilon,
1{+}\epsilon)A_t)\big]$ ($\epsilon{=}0.2$) is combined with a value-regression term
(only over steps where the agent chose to continue rather than terminate, when
\code{-use\_a\_t} is set) and an entropy bonus, both on separate coefficients
$\alpha_v, \alpha_h$ decayed on their own schedules independent of $\lambda_{\text{goal}}$.

% =====================================================================
\section{Live LLM-Explainability Monitoring}
\label{sec:explainability}
% =====================================================================

A per-epoch, always-on (not gated behind a verbosity flag) console report --
\emph{the LLM Explainability Brief} -- surfaces two categories of live diagnostic
signal that would otherwise only be visible by manually inspecting tensors:

\begin{enumerate}[leftmargin=1.4em]
  \item \textbf{Manager subgoal-selection histogram.} Every realized subgoal
    selection during the epoch's rollouts (\code{raw\_a[2]} from
    Algorithm~\ref{alg:forward}) is accumulated into an 8-bin histogram and printed
    as a percentage breakdown against the fixed ontology names -- but \emph{only}
    when \code{--llm\_distill} is active, since only then is the subgoal head being
    trained toward that specific name binding; without distillation, the same
    histogram is explicitly labeled as anonymous positional-index counts rather than
    silently mislabeled with names the policy was never trained to respect.
  \item \textbf{Distillation/shaping diagnostics.} When \code{--llm\_distill} is
    active: the current $\lambda_{\text{goal}}(t)$ (Eq.~\ref{eq:lambda-goal-sched}),
    the epoch's mean $\mathcal{L}_{\text{goal}}$ (or an explicit ``$\lambda_{\text{goal}}
    \approx 0$ this window'' note rather than a stale number), and the dataset size in
    use. When \code{--llm\_shaping} is active: the shaping weight $\lambda_t$ and the
    ratio $|F_t| / |r_{\text{env},t}|$ -- a direct, checkable signal for whether the
    shaping term is contributing a measurable gradient relative to the raw
    environment reward, rather than inferring a ``no measurable effect'' outcome only
    from final performance.
\end{enumerate}

\noindent Example output (this project's own training run, epoch 22 of 200):
\begin{lstlisting}[caption={Explainability brief, early training (11\% through the run)}]
[LLM EXPLAINABILITY BRIEF] Epoch 22
  Manager subgoal selections this epoch:
    DISCOVER_SUBNET           91.8%  (93971)
    ENUMERATE_HOST             6.9%  (7039)
    GAIN_INITIAL_ACCESS        1.0%  (1056)
    ESCALATE_PRIVILEGE         0.0%  (15)
    PIVOT                      0.1%  (97)
    CAPTURE_SENSITIVE_HOST     0.2%  (154)
    REDUCE_DETECTION           0.1%  (67)
    RECOVER_OR_REPLAN          0.0%  (1)
  llm_distill (manager-level): lambda_goal=0.9114  mean L_goal=0.4629  dataset=2000 records
\end{lstlisting}

\begin{lstlisting}[caption={Same run, epoch 200 (final) -- subgoal-selection collapse}]
[LLM EXPLAINABILITY BRIEF] Epoch 200
  Manager subgoal selections this epoch:
    DISCOVER_SUBNET          100.0%  (102400)
    ENUMERATE_HOST              0.0%  (0)
    GAIN_INITIAL_ACCESS         0.0%  (0)
    ESCALATE_PRIVILEGE          0.0%  (0)
    PIVOT                       0.0%  (0)
    CAPTURE_SENSITIVE_HOST      0.0%  (0)
    REDUCE_DETECTION            0.0%  (0)
    RECOVER_OR_REPLAN           0.0%  (0)
  llm_distill (manager-level): lambda_goal=0.0000  mean L_goal=n/a (lambda_goal ~ 0 this window)  dataset=2000 records
\end{lstlisting}

This is not a cosmetic feature: it is the mechanism that made the subgoal-selection
collapse analyzed in Section~\ref{sec:results-collapse} directly observable during
training, in human-readable form, epoch by epoch -- rather than requiring a
post-hoc statistical trawl through the raw JSONL step log to notice it. A curriculum
status banner (network realism-stage parameters: IDS enabled/decay/thresholds, scan
noise false-positive/negative rates, network-reliability timeout probability, service
churn probability) is printed alongside it at early epochs, every 10th epoch, and at
scenario-declared curriculum stage transitions, giving a complete picture of both
\emph{what the environment is doing to the agent} and \emph{what the agent's manager
is choosing to do} at every reported point in training.

% =====================================================================
\section{Experimental Setup}
\label{sec:setup}
% =====================================================================

\begin{table}[H]
\centering
\caption{Training run configuration (run id \texttt{b33undf9}).}
\label{tab:hparams}
\small
\begin{tabular}{@{}l p{9.7cm}@{}}
\toprule
Scenario (training) & \texttt{corp\_100hosts\_dynamic}\{base,\_varA,\_varB\}.v2.yaml (mixed) \\
Test scenario (in-loop \code{eval\_tst}) & \texttt{corp\_100hosts\_dynamic}.v2.yaml (base) \\
Architecture & \code{NASimNetDHRL} (212{,}640 params), \code{graph\_v2} obs., \code{-use\_a\_t} \\
Episode step limit & 400 \\
Parallel environments (\code{-batch}) & 128, over 8 CPU workers \\
PPO & $\gamma{=}0.99$, ppo\_k$=3$, ppo\_t$=8$, $\epsilon{=}0.2$, $\alpha_v{=}0.0333$ \\
Learning rate & $7\times 10^{-4}$, decayed $\times0.8$ every 10{,}000 steps, floor $3\times10^{-4}$ \\
Entropy coeff.\ $\alpha_h$ & $0.02$, time-decayed every 15{,}000 steps ($\times 0.5$), floor $0.005$ \\
Epoch length / max epochs & 100 steps/epoch, 200 epochs (20{,}000 train steps total) \\
LLM distillation & enabled; warm-start $E_w{=}3$ epochs; $\rho{=}8000$; $f{=}0.1$ (Stage 4 from step 18{,}000) \\
Teacher dataset & \code{training\_data/llm\_teacher\_dataset} (unsplit, all valid records) \\
Seed & 1 \\
Device & CPU (8 threads/process, BLAS pinned to 1 thread/process) \\
\bottomrule
\end{tabular}
\end{table}

\subsection{Teacher dataset}
The dataset used for distillation was collected via Algorithm~\ref{alg:collect}
against the single-model LLM backend (\code{qwen2.5:3b-instruct}). Table~\ref{tab:dataset-audit}
summarizes the audited dataset actually consumed by the training run in
Table~\ref{tab:hparams}.

\begin{table}[H]
\centering
\caption{Teacher dataset audit (\code{llm\_teacher/audit\_dataset.py}, this project's dataset).}
\label{tab:dataset-audit}
\begin{tabular}{lr}
\toprule
Total records & 2{,}235 \\
Valid & 2{,}000 (89.5\%) \\
Rejected & 235 (10.5\%) \\
\quad -- executable text in rationale & 124 \\
\quad -- horizon out of range & 105 \\
\quad -- API error (HTTP 500) & 6 \\
Distinct episodes & 170 \\
Train / test split (episode-level) & 2{,}000 / 235 records \\
Teacher backend/model & \code{qwen2.5:3b-instruct} (100\% of records) \\
Total labeling time & 17.1 min (0.46 s/record mean) \\
\bottomrule
\end{tabular}
\end{table}

\begin{table}[H]
\centering
\caption{Valid-record class balance and trigger mix -- both show the same
imbalance the collapse in Section~\ref{sec:results-collapse} tracks.}
\label{tab:class-balance}
\small
\begin{tabular}{lrl||lr}
\toprule
Goal class & Count & Mean conf. & Trigger & Count \\
\midrule
\code{DISCOVER\_SUBNET} & 1683 & 0.83 & \code{first\_compromise} & 805 \\
\code{ENUMERATE\_HOST} & 233 & 0.70 & \code{periodic} & 559 \\
\code{GAIN\_INITIAL\_ACCESS} & 55 & 0.94 & \code{high\_student\_entropy} & 439 \\
\code{ESCALATE\_PRIVILEGE} & 4 & 0.96 & \code{new\_subnet\_discovered} & 215 \\
\code{CAPTURE\_SENSITIVE\_HOST} & 12 & 0.92 & \code{episode\_start} & 170 \\
\code{REDUCE\_DETECTION} & 9 & 0.95 & \code{privilege\_change} & 39 \\
\code{PIVOT} & 4 & 0.84 & \code{three\_consecutive\_failures} & 8 \\
\code{RECOVER\_OR\_REPLAN} & 0 & n/a & & \\
\bottomrule
\end{tabular}
\end{table}

Five of eight classes fall below the audit tool's default minimum-class-count
threshold (20); \code{RECOVER\_OR\_REPLAN} has \emph{zero} valid labeled examples in
this dataset. This is flagged by \code{audit\_dataset.py} itself at collection time
(``training the distillation head against these classes now will mostly be
noise'') and is directly relevant to Section~\ref{sec:results-collapse}.

% =====================================================================
\section{Results}
\label{sec:results}
% =====================================================================

\subsection{Training dynamics}

\begin{table}[H]
\centering
\caption{Selected training-log snapshots (\texttt{training\_data/runs/b33undf9.json}).
\code{eval\_tst} is computed on the single base scenario every logged epoch.}
\label{tab:training-log}
\small
\begin{tabular}{@{}rrrrrr@{}}
\toprule
Train step & Loss & Entropy & \makecell[r]{eval\_tst\\reward\_avg\_ep.} & \makecell[r]{eval\_tst\\eplen\_avg} & \makecell[r]{eval\_tst\\captured\_avg} \\
\midrule
100    & 1.98  & 5.43 & --   & 400.0 & 24.8 \\
1000   & 8.84  & 3.95 & --   & 400.0 & 28.9 \\
14400  & 2.49  & 7.37 & 270.7 & 400.0 & 31.1 \\
19800  & --    & --   & --   & 115.2 & 29.6 \\
20000 (final) & 2.91 & 4.67 & 289.4 & 105.9 & 30.0 \\
\bottomrule
\end{tabular}
\end{table}

\subsection{Subgoal-selection collapse}
\label{sec:results-collapse}

The manager subgoal-selection histogram (Section~\ref{sec:explainability}) already
showed 91.8\% \code{DISCOVER\_SUBNET} at epoch 22 (11\% through training); by the
final epoch this reached \textbf{100.0\%}, with all seven other named goals --
including \code{CAPTURE\_SENSITIVE\_HOST}, the goal most directly tied to the task's
objective -- at exactly zero selections. This mirrors the teacher dataset's own class
imbalance almost exactly (84.2\% \code{DISCOVER\_SUBNET} in Table~\ref{tab:class-balance}):
a manager repeatedly distilled toward a heavily \code{DISCOVER\_SUBNET}-dominated
target, combined with the persistence \emph{switch} gate (Section~\ref{sec:arch})
having no explicit pressure against staying put, is a plausible, evidence-consistent
mechanism for the collapse, though this report does not run the random-label control
(condition C5, Section~\ref{sec:distill}) needed to establish it as the sole cause
rather than a contributing one. Despite the collapse in the \emph{labeled} subgoal
signal, the worker continued to make tangible task progress (\code{captured\_avg}
$\approx 30$ hosts throughout), consistent with the worker's action policy retaining
useful behavior largely independently of which subgoal label the manager was
nominally emitting at a given moment.

\subsection{Checkpoint-selection pitfall}
By default (\code{-save\_best\_metric captured\_avg}), the ``best'' checkpoint
(\code{model\_best.pt}) is saved whenever \code{eval\_tst.captured\_avg} sets a new
running maximum on the single, fixed, easiest training scenario. In this run, that
metric peaked at train step 14400 (Table~\ref{tab:training-log}) and was never
strictly exceeded again -- even though \code{eval\_tst.eplen\_avg} dropped from
$400.0$ (every episode exhausting the step budget) to $105.9$ by the final step,
meaning the policy had learned to finish episodes roughly $3.8\times$ faster for a
comparable capture count. \code{captured\_avg} alone cannot see an episode-length
improvement, so \code{model\_best.pt} silently froze at a materially worse-behaving
policy than the run's own final checkpoint, \code{model.pt}.

\subsection{Cross-scenario generalization}
Both checkpoints were evaluated on all five scenario instances
(Section~\ref{sec:env}), 100 held-out episodes per scenario, greedy (\code{net.eval()})
action selection.

\begin{table}[H]
\centering
\caption{\code{model\_best.pt} (frozen at train step 14400) vs.\ \code{model.pt}
(final, train step 20000). \code{reward\_avg\_episodes} is the sum of per-episode
reward, not the (misleading, unequal-episode-length) per-step average.}
\label{tab:eval-comparison}
\begin{tabular}{lrrrrrr}
\toprule
& \multicolumn{3}{c}{\code{model\_best.pt}} & \multicolumn{3}{c}{\code{model.pt} (final)} \\
\cmidrule(lr){2-4} \cmidrule(lr){5-7}
Scenario & reward\_avg\_ep. & eplen & captured & reward\_avg\_ep. & eplen & captured \\
\midrule
base (train)    & 3.24   & 400.0 & 30.2 & \textbf{259.1} & \textbf{111.6} & 29.6 \\
varA (train)    & 21.8   & 400.0 & 17.5 & \textbf{80.5}  & \textbf{271.2} & 17.5 \\
varB (train)    & $-$12.2 & 400.0 & 14.4 & \textbf{46.0}  & \textbf{272.7} & 13.3 \\
test (held-out) & $-$36.7 & 399.4 & 11.9 & \textbf{$-$8.3} & \textbf{277.8} & 9.2 \\
bridge (held-out) & $-$22.1 & 400.0 & 25.9 & \textbf{205.5} & \textbf{117.9} & 24.7 \\
\bottomrule
\end{tabular}
\end{table}

\code{captured\_avg} is essentially flat between the two checkpoints on every
scenario (within evaluation noise over 100 episodes) -- the final checkpoint is
\emph{not} trading capture ability for speed, it is achieving the same capture
ability in far fewer wasted steps. \code{reward\_avg\_episodes} flips from negative
to strongly positive on 3 of 5 scenarios and improves substantially on the remaining
two (including the held-out \code{test} scenario, still the weakest under both
checkpoints, consistent with it having the lowest mean exploit probability in
Table~\ref{tab:scenario-difficulty}). The practical conclusion is direct: for this
architecture and this choice of \code{-save\_best\_metric}, the plain final
checkpoint generalizes measurably better across held-out scenarios than the
metric-selected ``best'' one, and \code{reward\_avg\_episodes} (which penalizes an
episode that never terminates) would be the more appropriate optimization target for
\code{-save\_best\_metric} in future runs of this system.

\subsection{Comparison against a non-distilled Pen-DHRL baseline}
\label{sec:baseline-comparison}
Beyond the checkpoint-selection question above, \pendhrl{}'s final checkpoint was
compared against a plain \code{NASimNetDHRL} baseline trained with \emph{no} LLM
distillation -- identical architecture and parameter count (212k), isolating the
distillation intervention itself. The baseline was trained twice independently; the
value reported per scenario is from its stronger run, so it is shown in its best
light rather than penalized by an unlucky run. All four evaluation metrics are
reported for both, not just \code{reward\_avg\_episodes}. No baseline run was
collected for the held-out \code{test} scenario, so it is omitted here (Pen-DHRL-TD
alone on \code{test} is already in Table~\ref{tab:eval-comparison} above).

\begin{table}[H]
\centering
\caption{\pendhrl{} (final checkpoint).}
\label{tab:baseline-comparison-td}
\small
\begin{tabular}{lrrrr}
\toprule
Scenario & reward\_avg & reward\_avg\_episodes & eplen\_avg & captured\_avg \\
\midrule
base   & 2.3216 & 259.12 & 111.61 & 29.63 \\
varA   & 0.2967 & 80.47  & 271.18 & 17.50 \\
varB   & 0.1688 & 46.04  & 272.69 & 13.25 \\
bridge & 1.7425 & 205.48 & 117.93 & 24.67 \\
\bottomrule
\end{tabular}
\end{table}

\begin{table}[H]
\centering
\caption{Pen-DHRL baseline -- same architecture, same parameter count, no
\code{--llm\_distill}.}
\label{tab:baseline-comparison-base}
\small
\begin{tabular}{lrrrr}
\toprule
Scenario & reward\_avg & reward\_avg\_episodes & eplen\_avg & captured\_avg \\
\midrule
base   & 0.5428 & 217.11 & 400.0 & 29.59 \\
varA   & 0.1708 & 68.33  & 400.0 & 13.73 \\
varB   & 0.0778 & 31.12  & 400.0 & 9.39 \\
bridge & 1.1470 & 161.96 & 141.20 & 23.36 \\
\bottomrule
\end{tabular}
\end{table}

\begin{figure}[H]
\centering
\includegraphics[width=0.85\textwidth]{eval_results_model_best.png}
\caption{Both tables above, rendered directly from the evaluation output.}
\label{fig:eval-comparison}
\end{figure}

\pendhrl{} wins on every metric (\code{reward\_avg}, \code{reward\_avg\_episodes},
\code{captured\_avg}) on every scenario, and terminates well before the 400-step
limit on all four, while the baseline only does so on \code{bridge}. Because the two
runs share architecture and parameter count exactly, this isolates the LLM-teacher
distillation signal itself as the source of the gain, rather than any difference in
model capacity.

% =====================================================================
\section{Discussion and Limitations}
% =====================================================================

\begin{itemize}[leftmargin=1.4em]
  \item \textbf{Class imbalance in the teacher dataset.} Five of eight ontology
    classes are below a reasonable minimum count for reliable supervision, and
    \code{RECOVER\_OR\_REPLAN} has zero examples (Table~\ref{tab:class-balance}).
    The trigger set (Algorithm~\ref{alg:collect}) over-samples early-episode/
    exploration-phase states relative to rare escalation/pivot moments; a targeted
    collection pass using \code{--policy checkpoint} against a partially-trained
    agent that has already learned to reach those rarer states (rather than the
    \code{random}/\code{untrained\_net} policies that dominate easy states) is the
    documented mitigation path.
  \item \textbf{Distillation-collapse hypothesis, not proof.} Section~\ref{sec:results-collapse}'s
    causal link between dataset imbalance and the subgoal-selection collapse is
    plausible and evidence-consistent but not isolated from the persistence-gate
    dynamics or the RL objective's own contribution. The random-label control
    (\code{--llm\_distill\_shuffle\_labels}, condition C5) and the heuristic-teacher
    ablation (Section~\ref{sec:teacher}) both already exist in this codebase to
    disambiguate this and were not run for the present report.
  \item \textbf{Single teacher model.} All reported results use one 3B-parameter
    local model. The cascade escalation path (Section~\ref{sec:teacher}) and
    workstation-scale larger-model candidates are implemented but unevaluated here.
  \item \textbf{Potential-based shaping unused.} Section~\ref{sec:shaping}'s reward
    level channel was not enabled for this run; its interaction with the collapse
    observed in distillation-only training is unexplored.
  \item \textbf{Checkpoint-selection metric.} As shown directly in
    Table~\ref{tab:eval-comparison}, \code{-save\_best\_metric captured\_avg} is not
    a safe default for this task; it is blind to episode-length degeneracy.
  \item \textbf{Baseline comparison is one run.} Table~\ref{tab:baseline-comparison-base}
    reports the non-distilled baseline's stronger of two independently trained runs,
    not a distribution over many runs -- a useful signal, not a statistically
    rigorous claim of significance. The same caveat applies symmetrically to
    \pendhrl{} itself: the results throughout this report come from a single
    training run.
\end{itemize}

% =====================================================================
\section{Conclusion}
% =====================================================================

\pendhrl{} demonstrates a complete, offline-only pipeline for grounding a
hierarchical RL agent's otherwise-anonymous subgoal space in LLM-derived semantic
supervision: a fixed goal ontology, a sanitized/validated/provenance-tracked teacher
dataset built via triggered sampling under a diversified policy mixture, a
confidence-weighted distillation loss folded into a four-stage training schedule
alongside standard PPO, and a live explainability mechanism that surfaces the
manager's realized behavior every epoch rather than only at evaluation time. Applied
to a 100-host dynamic corporate scenario family, this pipeline surfaced concrete,
actionable findings that would not have been visible from aggregate reward curves
alone: a subgoal-selection collapse traceable to the teacher dataset's own class
imbalance, and a checkpoint-selection metric blind to episode-length degeneracy that
caused the saved ``best'' model to substantially underperform the run's own final
policy on held-out generalization. Set against an independently trained,
same-architecture, same-parameter-count Pen-DHRL baseline run without the teacher
signal, the final distilled policy also wins on every reported metric on every
scenario tested -- evidence that the gain traces to the distillation signal itself
rather than model capacity. All of these findings, and the tooling that surfaced them,
are directly actionable
for the next data-collection and training pass.

\end{document}
